\begin{table}[H]
\caption{HW Interfaces of Module InstrDec0}
\label{tab:asDecPer02}
\centering
\begin{tabularx}{\textwidth}{|l |l |X|}
  \hline
   & Module internal name & Chip level name \\
  \hline
  \hline
  Supply &  &  \\
  \hline
  Clock & - & - (pure combinational block, no \texttt{clk\_i}/\texttt{rst\_i}) \\
  \hline
  Bus & - & - \\
  \hline
  Interrupt & \texttt{trap\_illegal\_instr\_o} &
    \textit{RV\_NoPipeline} (\texttt{asCPUx.sv}): \texttt{trap\_illegal\_instrx\_s}.
    \textit{RV\_PipelineCacheFwd} (\texttt{cc\_cpupipefwd.sv}):
    \texttt{id\_trap\_illegal\_s}. Both feed \texttt{mcause}\,=\,2 directly;
    see Section~\ref{subsec:irqconcept}. \\
  \hline
  DMA & - & - \\
  \hline
  External Signals & - & - \\
  \hline
  Others & \texttt{instr\_opcode\_i}(6:0) &
    \textit{NoPipeline}: \texttt{ir\_s}(6:0). \textit{PipelineCacheFwd}:
    \texttt{if\_id\_instr\_r}(6:0) \\
         & \texttt{instr\_func3\_i}(2:0) &
    \texttt{ir\_s}(14:12) / \texttt{if\_id\_instr\_r}(14:12) \\
         & \texttt{instr\_func7b5\_i} &
    \texttt{ir\_s}(30) / \texttt{if\_id\_instr\_r}(30) \\
         & \texttt{take\_i} &
    \textit{NoPipeline}: \texttt{take\_s} (branch-taken feedback from the
    branch ALU). \textit{PipelineCacheFwd}: tied to \texttt{1'b0}
    (unused -- kept only for interface compatibility, see the
    \texttt{verilator lint\_off UNUSEDSIGNAL} note in
    \texttt{as\_instr\_decode.sv}) \\
         & \texttt{mux\_resultSrc\_o} (\texttt{result\_src\_t}, 2\;bit) &
    \textit{NoPipeline}: \texttt{resultSrcx\_s}. \textit{PipelineCacheFwd}:
    \texttt{id\_resultSrc\_s} \\
         & \texttt{en\_dMemWr\_o} & \texttt{dMemWr\_s} / \texttt{id\_dMemWr\_s} \\
         & \texttt{en\_dMemRd\_o} & \texttt{dMemRd\_s} / \texttt{id\_dMemRd\_s} \\
         & \texttt{mux\_aluSrcB\_o} & \texttt{aluSrcB\_s} / \texttt{id\_aluSrcB\_s} \\
         & \texttt{mux\_aluSrcA\_o} (\texttt{mux\_a\_t}, 2\;bit) &
    \texttt{aluSrcA\_s} / \texttt{id\_aluSrcA\_s} \\
         & \texttt{en\_regWr\_o} & \texttt{regWr\_s} / \texttt{id\_regWr\_s} \\
         & \texttt{mux\_jump\_o} & \texttt{jump\_s} / \texttt{id\_jump\_s} \\
         & \texttt{sel\_immSrc\_o} (\texttt{imm\_src\_t}, 3\;bit) &
    \texttt{immSrcx\_s} / \texttt{id\_immSrc\_s} \\
         & \texttt{alu\_op\_o} (\texttt{alu\_op\_t}, 6\;bit declared,
           15 values used) & \texttt{aluSela\_s} / \texttt{id\_aluOp\_s} \\
         & \texttt{br\_op\_o} (\texttt{br\_op\_t}, 3\;bit) &
    \texttt{aluSelb\_s} / \texttt{id\_brOp\_s} \\
  \hline
\end{tabularx}
\end{table}
\noindent\footnotesize Both chip variants instantiate the identical,
shared \texttt{as\_instr\_decode} module (\texttt{RWU\_RV/RV\_NoPipeline/
src/asInstrDecode.sv}); only the chip-level net names at the instantiation
point differ, as listed above.\normalsize
